\documentclass{article}

% --- NeurIPS style if available, else a clean article fallback so it always
% --- compiles. Drop neurips_2024.sty next to this file for the real format.
\IfFileExists{neurips_2024.sty}{%
  \usepackage[preprint]{neurips_2024}% loads natbib internally
}{%
  \usepackage[margin=1in]{geometry}%
  \usepackage{times}%
  \usepackage{natbib}%
  \newcommand{\AND}{\quad}%
}

\usepackage[utf8]{inputenc}
\usepackage[T1]{fontenc}
\usepackage{hyperref}
\usepackage{url}
\usepackage{booktabs}
\usepackage{amsmath}
\usepackage{amssymb}
\usepackage{graphicx}
\usepackage{xcolor}
\graphicspath{{figures/}{../results/}}

% Include a figure if present, else draw a labeled placeholder so the document
% always compiles (copy results/*.png into paper/figures/ for the real figures).
\newcommand{\figorbox}[1]{%
  \IfFileExists{figures/#1}{\includegraphics[width=\linewidth]{#1}}{%
  \IfFileExists{../results/#1}{\includegraphics[width=\linewidth]{#1}}{%
  \fbox{\parbox[c][0.62\linewidth][c]{0.95\linewidth}{\centering\ttfamily #1\\(run pipeline + copy to figures/)}}}}%
}

\title{How Far Can Cheap Linear Transforms Close the\\
Human--Model Visual Similarity Gap?}

\author{%
  Sudarssan N. \\
  \texttt{sudarssan73@gmail.com} \\
}

\begin{document}
\maketitle

\begin{abstract}
Pretrained vision models encode object similarity differently from humans. We
ask how much of this gap is a simple linear distortion of an otherwise
human-like space, and whether the correction is shared across architectures.
Using the THINGS odd-one-out benchmark (4.70M human triplet judgments over 1{,}854
objects), we evaluate five frozen backbones spanning self-supervised (DINOv2),
contrastive (CLIP, SigLIP), and supervised (ViT, ResNet-50) training, and learn a
single L2-regularized linear transform per model with a triplet objective. We
find three things. (i)~\textbf{Universal near-ceiling alignability}: a cheap
linear transform lifts \emph{every} backbone to 88--92\% of the human noise
ceiling (0.59--0.62 accuracy), erasing large initial differences between training
paradigms; the largest gains accrue to the weakest starting points (DINOv2,
$+0.20$). (ii)~\textbf{Model-specificity}: transforms transfer poorly across
architectures---a transform learned on one model recovers only $\sim$15--25\% of
the within-model gain on another---indicating the gap is mostly
architecture/objective-specific rather than a single universal skew.
(iii)~\textbf{A shared semantic ceiling}: the same abstract/functional categories
(office supplies, furniture, musical instruments, medical equipment) resist linear
alignment across all backbones. Results are stable across seeds, regularization
strength, and a stricter image-disjoint split.
\end{abstract}

\section{Introduction}
Human alignment of internal representations is increasingly viewed as a driver of
robustness, interpretability, and transfer in vision models
\citep{sucholutsky2023alignment,muttenthaler2024aligning}. The THINGS initiative
\citep{hebart2019things,hebart2023things} provides large-scale behavioral object
similarity data through a triplet odd-one-out task, operationalizing human object
similarity at scale. We use it to study a focused, practically relevant question:
given \emph{frozen} embeddings from standard pretrained backbones, how far can a
\emph{single cheap linear transform} push human--model alignment, and does such a
transform transfer across architectures? A positive transfer result would imply
the misalignment is largely a shared linear distortion; a negative one would
imply deeper, model-specific structure.

\paragraph{Contributions.}
(1)~A reproducible THINGS odd-one-out pipeline over a five-model zoo, fully
scripted from public data. (2)~Evidence that a cheap linear transform brings all
backbones to within 88--92\% of the human ceiling, with the biggest gains for the
weakest models. (3)~A cross-model transfer analysis showing the alignment
transform is mostly model-specific. (4)~A category-level error analysis revealing
a shared semantic ceiling on linear alignment.

\section{Related work}
\citet{muttenthaler2023human,muttenthaler2024aligning} show that linear
transforms optimized against human similarity improve correspondence and
downstream robustness; we complement this by directly contrasting architectures
and quantifying cross-model transfer of the learned transform. The human
similarity structure we target derives from the SPoSE/THINGS embedding
\citep{hebart2020revealing,hebart2023things}. We compare representations using
odd-one-out accuracy and representational similarity analysis
\citep{kriegeskorte2008rsa}.

\section{Method}
\paragraph{Data.} We use the 1{,}854 THINGS reference objects with one
copyright-free (CC0) image per concept \citep{stoinski2024thingsplus}, and the
THINGS odd-one-out triplets \citep{hebart2023things}: 4.12M train, 0.45M
validation, and a 15.6k noise-ceiling test set. Triplets are 0-indexed with the
odd-one-out in the last position; the empirical human noise ceiling is $0.673$
and chance is $1/3$.

\paragraph{Backbones.} Five frozen models: DINOv2 ViT-B/14 (self-supervised),
CLIP ViT-B/16 and SigLIP ViT-B/16 (contrastive), ViT-B/16 and ResNet-50
(ImageNet-supervised) \citep{oquab2024dinov2,radford2021clip,zhai2023siglip,dosovitskiy2021vit,he2016resnet}.
We extract one embedding per image (CLS/pooled features).

\paragraph{Odd-one-out scoring.} For a triplet with chosen-similar pair $(a,b)$
and odd-one-out $o$, we score pairwise cosine similarities and predict the odd
item as the one least similar to the other two; accuracy is agreement with the
human choice.

\paragraph{Linear alignment.} We learn a transform $W$ on standardized features
with the SPoSE/VICE-style objective: with $z=Wx$ normalized, a 3-way softmax over
pair similarities $\{s_{ab},s_{ao},s_{bo}\}$ (scaled by a learnable temperature)
is trained by cross-entropy to prefer the human pair, with L2 regularization on
$W$. We early-stop on validation accuracy.

\paragraph{Cross-model transfer.} To compare transforms across differing
dimensionalities, we standardize and PCA-reduce every model to a shared 256-d
space, learn $W$ per model there, and evaluate accuracy when applying $W_{\text{src}}$
to $\text{model}_{\text{tgt}}$'s features.

\paragraph{RSA.} We correlate (Spearman) each model's $1854\times1854$ cosine
similarity matrix with the human SPoSE similarity matrix, before and after
alignment.

\section{Results}

\paragraph{A cheap transform reaches $\sim$90\% of the ceiling for all models.}
Table~\ref{tab:main} shows zero-shot vs.\ aligned odd-one-out accuracy. Despite
zero-shot accuracy ranging from 0.41 (DINOv2) to 0.47 (CLIP), all models converge
to 0.59--0.62 after a single linear transform (88--92\% of the 0.673 ceiling).
The largest improvements go to the weakest baselines, so initial differences
between training paradigms are largely a linear artifact.

\begin{table}[t]
\centering
\caption{Zero-shot vs.\ linearly-aligned odd-one-out accuracy on the THINGS
test set (human ceiling $0.673$, chance $0.333$). $\Delta$ is aligned minus
zero-shot; \% ceiling is aligned/ceiling.}
\label{tab:main}
\begin{tabular}{llccccc}
\toprule
Model & Family & $d$ & Zero-shot & Aligned & $\Delta$ & \% ceiling \\
\midrule
SigLIP ViT-B/16 & contrastive & 768 & 0.468 & \textbf{0.616} & $+0.148$ & 91.6 \\
ResNet-50       & supervised  & 2048& 0.433 & 0.613 & $+0.180$ & 91.1 \\
DINOv2 ViT-B/14 & self-sup.   & 768 & 0.408 & 0.609 & $\mathbf{+0.200}$ & 90.5 \\
CLIP ViT-B/16   & contrastive & 512 & \textbf{0.474} & 0.604 & $+0.130$ & 89.8 \\
ViT-B/16        & supervised  & 768 & 0.436 & 0.593 & $+0.157$ & 88.0 \\
\bottomrule
\end{tabular}
\end{table}

\paragraph{Transforms are mostly model-specific.}
Table~\ref{tab:transfer} reports the transfer matrix in the shared 256-d space.
Own-model transforms (diagonal) reach 0.55--0.60, but transferred transforms
(off-diagonal) reach only 0.39--0.47. While off-diagonal entries exceed the
identity baseline (0.35--0.39), they recover only $\sim$15--25\% of the
within-model alignment gain, with CLIP the strongest donor. The human--model gap
is therefore largely architecture/objective-specific, with only a weak shared
linear component (Fig.~\ref{fig:heatmap}).

\begin{table}[t]
\centering
\caption{Cross-model transfer (test accuracy, shared 256-d). Rows: source
transform $W$; columns: target features. Last row: identity baseline ($W{=}I$).}
\label{tab:transfer}
\begin{tabular}{lccccc}
\toprule
$W$ \textbackslash{} target & CLIP & DINOv2 & ResNet & SigLIP & ViT \\
\midrule
CLIP    & \textbf{0.593} & 0.412 & 0.440 & 0.417 & 0.424 \\
DINOv2  & 0.415 & \textbf{0.576} & 0.396 & 0.391 & 0.414 \\
ResNet  & 0.448 & 0.393 & \textbf{0.566} & 0.402 & 0.396 \\
SigLIP  & 0.431 & 0.398 & 0.393 & \textbf{0.604} & 0.389 \\
ViT     & 0.466 & 0.390 & 0.413 & 0.409 & \textbf{0.551} \\
\midrule
\textit{baseline} & 0.390 & 0.361 & 0.353 & 0.375 & 0.379 \\
\bottomrule
\end{tabular}
\end{table}

\paragraph{RSA converges after alignment.}
Raw Spearman RSA to the human embedding varies widely (DINOv2 0.14, ViT 0.15 vs.\
CLIP 0.42), but after alignment all models converge to 0.78--0.83
(Fig.~\ref{fig:rsa}). DINOv2 and supervised ViT are the most ``rotated away'' from
human geometry yet are recoverable with a linear map.

\begin{figure}[t]
\centering
\begin{minipage}{0.48\textwidth}
\centering
\figorbox{transfer_heatmap.png}
\caption{Cross-model transfer of alignment transforms (odd-one-out test acc).}
\label{fig:heatmap}
\end{minipage}\hfill
\begin{minipage}{0.48\textwidth}
\centering
\figorbox{rsa.png}
\caption{RSA to the human SPoSE embedding, before vs.\ after alignment.}
\label{fig:rsa}
\end{minipage}
\end{figure}

\paragraph{A shared semantic ceiling.}
Grouping residual errors by the THINGS 27 high-level categories, the hardest
categories after alignment are consistent across all five backbones:
\emph{office supply, furniture, kitchen tool, musical instrument, medical
equipment}, and \emph{toy}. Concrete categories (animals, food, fruit) are nearly
solved, while abstract/functional concepts resist linear alignment regardless of
architecture.

\paragraph{Robustness.} Aligned accuracies are stable across random seeds
(std $\le 0.004$; Table~\ref{tab:robust}) and insensitive to the L2 strength over
$\{10^{-4},10^{-3},10^{-2}\}$ (variation $<0.02$, optimum at $10^{-3}$ for all
models). To quantify image-level leakage we compare a strict \emph{image-disjoint}
split (train/validation/test share no images; $\sim$27\% of triplets form the
test set) against a \emph{size-matched random triplet split} of the same pool, so
the two differ only in whether train and test share images. The leakage gap is
small (Table~\ref{tab:robust}): $+0.013$ to $+0.056$, i.e.\ aligned accuracy
includes only a 1--6 point inflation from shared images. Even fully
image-disjoint, aligned accuracy (0.55--0.58) far exceeds the zero-shot baselines
(0.41--0.47), confirming the transform generalizes to unseen images.
Interestingly, the contrastive models leak least (CLIP $+0.013$, SigLIP $+0.020$)
while supervised/self-supervised models leak more (ResNet $+0.056$, DINOv2
$+0.035$, ViT $+0.032$), suggesting contrastive alignment is more concept-general
and less reliant on image-specific structure.

\begin{table}[t]
\centering
\caption{Robustness of aligned odd-one-out accuracy: stability across seeds, and
an image-leakage estimate from an image-disjoint vs.\ size-matched triplet-level
split of the same pool (gap = matched $-$ image-disjoint).}
\label{tab:robust}
\begin{tabular}{lcccc}
\toprule
Model & Seeds (mean$\pm$std) & Image-disjoint & Matched-triplet & Leakage gap \\
\midrule
CLIP ViT-B/16   & $0.605\pm0.001$ & 0.580 & 0.593 & $+0.013$ \\
SigLIP ViT-B/16 & $0.616\pm0.003$ & 0.582 & 0.602 & $+0.020$ \\
ViT-B/16        & $0.594\pm0.003$ & 0.551 & 0.583 & $+0.032$ \\
DINOv2 ViT-B/14 & $0.602\pm0.002$ & 0.555 & 0.590 & $+0.035$ \\
ResNet-50       & $0.619\pm0.004$ & 0.549 & 0.606 & $+0.056$ \\
\bottomrule
\end{tabular}
\end{table}

\section{Discussion and limitations}
A single cheap linear transform recovers most of the human--model similarity gap
for every backbone we tested, but the required transform is largely
model-specific and a shared semantic residual remains. This favors using
inexpensive per-model alignment in practice while cautioning against a
``universal'' alignment patch. Limitations: a single reference image per concept
(visual variation is not modeled), a behavioral---not neural---alignment target,
base-size models only, and linear maps only (nonlinear or hyperbolic maps may
behave differently). Extensions include THINGS-fMRI/MEG neural alignment and
geometry beyond Euclidean.

\section{Conclusion}
On THINGS odd-one-out, cheap linear transforms close most of the human--model
gap uniformly across architectures, but do not transfer well between them, and
leave a shared semantic ceiling---suggesting the residual misalignment is
structural rather than a simple universal distortion.

\paragraph{Reproducibility.} Code and scripts (data download, feature
extraction, alignment, transfer, analysis, robustness) are available at
\url{https://github.com/Sudarssan-N/VisionModal-HumalObjectSimilarity}.

\bibliographystyle{plainnat}
\bibliography{references}

\end{document}
